\section{Introduction}
\label{sec:intro}

% Robotic manipulation is one of the hardest problems in embodied intelligence precisely because it demands that perception, language understanding, and physical control operate as a unified system rather than independent modules. Vision-Language-Action models have emerged as the natural vehicle for this integration, inheriting the semantic and perceptual depth of modern VLMs and extending them into the action domain. The implicit promise of a VLA model is that of any foundation model: train broadly, transfer efficiently, generalize reliably.


Foundation models in language and multimodality~\citep{gpt3,gpt4,dubey2024llama,gemini,yang2025qwen3,qwen3-vl,qwen35blog} achieve strong generalization because heterogeneous data sources can be aligned under a unified formulation, and abundant low-cost data from the internet allows diverse training signals to reinforce one another at scale. 
In this report, we investigate whether this scaling recipe can be applied to robotic manipulation to achieve genuine generalization. 
This is challenging because, unlike text, manipulation data is heterogeneous by nature, expensive to collect, and narrow in diversity, making alignment and scale simultaneously difficult to achieve.

% That promise remains largely unfulfilled. Despite a rapid succession of VLA models reporting competitive numbers on standard manipulation benchmarks, the generalization being demonstrated is shallow. Out-of-distribution evaluations in most published work amount to swapping object instances or backgrounds while keeping the same embodiment, task structure, and workspace layout as training. More tellingly, on several standard benchmarks, fine-tuning a pretrained VLA on the provided training split yields performance comparable to training from scratch on the same data~\citep{Yan_2025_ICCV}. A model whose pretrained prior offers no advantage over random initialization at transfer time cannot meaningfully be called a foundation model. What practitioners consistently find most useful in real deployment is not benchmark rank but how efficiently a model adapts to a new task with modest domain-specific data. The correct criterion for evaluating a VLA foundation model is therefore not in-distribution task success but out-of-distribution transfer efficiency and robustness.



The current state of Vision-Language-Action (VLA) models~\citep{openvla,bjorck2025grootn1,black2024pi0,pi05,community2026starvla,ddpvla,kim2025fine} illustrates how far this gap remains. 
Despite a rapid succession of models reporting competitive numbers on standard benchmarks~\citep{liu2023libero,nasiriany2024robocasa,mu2025robotwin}, the generalization being demonstrated is largely superficial~\citep{zhang2026world}. 
OOD evaluations in most works involve only minor visual perturbations while preserving the same embodiment, task structure, and workspace layout as data collection, and performance degrades sharply when models are tested beyond these narrow settings.
%Out-of-distribution evaluations in most published work amount to swapping object instances or backgrounds while preserving the same embodiment, task structure, and workspace layout as training. 
%When these models encounter genuinely novel conditions, such as unseen embodiments, novel camera viewpoints, cluttered scenes, or compositionally new instructions, performance degrades sharply, revealing that the pretrained priors have not internalized transferable manipulation structure.
%Moreover, fine-tuning a pretrained VLA on a benchmark's own training split yields performance comparable to training from scratch on the same data~\citep{Yan_2025_ICCV,community2026starvla}, indicating that the pretrained prior contributes negligible transferable value under in-domain evaluation.


The reason these pretrained priors fail to transfer is twofold. 
First, existing robotic demonstration corpora~\citep{padalkar2023openxembodiment,khazatsky2024droid,fang2023rh20t} are concentrated in narrow teleoperation setups, far too limited in embodiment and task diversity for a scaling recipe to take hold. Second, and more fundamentally, data diversity alone is insufficient without alignment~\citep{luo2026joint, wang2026rethinking}.
When demonstrations from different embodiments arrive with incompatible observation and action representations, scaling data volume produces interference rather than synergy. Prior cross-embodiment efforts~\citep{bjorck2025grootn1,zheng2025xvla,pi05} have adopted shared architectures, embodiment tokens, or unified action tokenization, but without a formulation that makes the same physical motion numerically consistent across embodiments, additional data cannot be converted into improved capability.
Alignment is therefore not an independent engineering choice but a prerequisite for data scaling itself.

%The underlying reason is straightforward. Existing robotic demonstration collections are concentrated in narrow, costly teleoperation setups, and no prior training formulation has been designed to absorb demonstrations from fundamentally incompatible action and state spaces in a way that produces coherent cross-embodiment representations. Without both, the scaling recipe cannot take hold.

% Understanding why existing VLA models fall short of this criterion points directly to what must be fixed. The fundamental obstacle is data. Genuine generalization in any domain requires training on data diverse enough to force the model to learn structure rather than memorize patterns. In robotic manipulation, achieving this diversity means ingesting demonstrations across embodiments, camera configurations, action spaces, and task contexts simultaneously. Without training data that spans this diversity at scale, transferability remains fundamentally limited. The challenge is compounded by the fact that robotic demonstration collection is expensive and concentrated in narrow setups, leaving the vast reservoir of egocentric human manipulation video largely untapped.

% Scaling across heterogeneous data sources in turn demands alignment. A model trained on demonstrations from fundamentally incompatible action and state spaces will not naturally learn cross-embodiment structure without explicit design choices that make diverse data mutually reinforcing rather than mutually interfering. Representation, motion, and behavioral alignment are not independent engineering decisions but interconnected necessities for multi-source training to produce transferable representations.


We present \textbf{\ours}, a Vision-Language-Action foundation model built upon Qwen-VL~\citep{yang2025qwen3, qwen3-vl,qwen35blog}, designed around this principle: \emph{alignment first, then scale}. %designed from the ground up around the goal of genuine generalizability. 
\ours introduces a unified alignment framework across the representation, motion, and behavioral dimensions of manipulation, making large-scale multi-source training coherent rather than conflicting: a canonical state-action representation with per-dimension binary masking accommodates diverse robot morphologies within a single template, a camera-frame delta pose parameterization makes visually similar motions numerically proximate across coordinate frames, and an in-context policy adaptation mechanism reads intra-episode execution history as an implicit embodiment identifier for kinematic-aware behavioral adjustment. 
% （先讲alignment, data scaling放到下一段？）
%To build a data corpus with the breadth that generalization demands, \ours combines a human-to-robot synthesis pipeline that renders egocentric hand demonstrations as robot trajectories across 15 platforms with a rigorous multi-stage curation pipeline that harmonizes heterogeneous real-robot and simulation datasets, accumulating $\sim$30,000 hours of manipulation data. 
%To extract consistent cross-embodiment signal from this diversity, \ours introduces a unified alignment framework spanning a canonical state-action representation with per-dimension binary masking, a camera-frame delta pose parameterization that makes visually similar motions numerically proximate across coordinate frames, and an in-context policy adaptation mechanism that reads intra-episode execution history as an implicit embodiment identifier for kinematic-aware behavioral adjustment. 
Training is conducted under a dual-stream co-training strategy that jointly optimizes over manipulation data and a curated vision-language stream, preventing the VLM backbone's perceptual and reasoning capabilities from eroding under action prediction pressure.

This alignment capability in turn enables \ours to absorb manipulation data at a scale that prior training regimes could not sustain. 
A human-to-robot synthesis pipeline converts egocentric hand demonstrations into robot trajectories across 15 platforms, and a rigorous curation pipeline harmonizes heterogeneous real-robot and synthetic datasets, together accumulating $\sim$38,100 hours of manipulation data. 
Notably, this entire corpus is constructed from \textbf{only} open-source robotic manipulation datasets and egocentric human videos without any proprietary data collection, yet \ours already exhibits emergent generalization capabilities, including robustness to perturbations, zero-shot instruction following, reactive error recovery, and cross-embodiment transfer.

We further argue that the field's evaluation methodology must evolve alongside its models. Standard in-domain benchmarks, where models without large-scale robot pretraining match or exceed pretrained ones~\citep{Yan_2025_ICCV,community2026starvla}, systematically fail to distinguish genuine generalization from in-distribution pattern memorization.
%We also argue that the field needs better tools for measuring what it claims to care about. In-domain benchmark performance is a poor proxy for the generalization capability that makes a foundation model useful. 
\ours therefore introduces OOD evaluation settings including LIBERO-Plus~\citep{fei2025liberoplus}, RoboTwin-Clean2Rand~\citep{mu2025robotwin}, RoboCasa365~\citep{nasiriany2026robocasa365}, and EBench~\citep{ebench2026}, alongside two new benchmarks: \textbf{RoboTwin-IF}, an instruction-following benchmark that tests whether policies condition on language as a genuine control signal rather than exploit visual shortcuts, and \textbf{RoboTwin-XE}, which evaluates zero-shot transfer to morphologically distinct robots. 
\ours achieves state-of-the-art on standard benchmarks and substantially outperforms existing VLA models including GR00T-N1.7 and $\pi_{0.5}$ across all OOD evaluation axes, ranks 1st on the RoboChallenge Table30-v1 generalist track with a 20\% relative improvement, and is validated on various real-robot platforms and tasks. We hope these results and benchmarks together raise the standard for how VLA generalization is measured across the field.

Our contributions are as follows.

\textbf{A unified alignment framework for cross-embodiment training.} We address representational heterogeneity through three complementary mechanisms. A canonical state-action representation with per-dimension binary masking, a camera-frame delta pose parameterization that grounds end-effector actions in the visual domain, and an in-context policy adaptation mechanism that treats intra-episode execution history as an implicit embodiment identifier together enable consistent signal extraction across diverse embodiments. A dual-stream co-training strategy jointly optimizes manipulation and vision-language objectives to preserve the perceptual and reasoning capabilities that underpin generalization.

\textbf{A scalable multi-source data corpus.} We consolidate $\sim$38,100 hours of manipulation data from open-source robot datasets and egocentric human demonstrations. A human-to-robot synthesis pipeline converts any egocentric demonstration into trajectories across 15 robot platforms, providing a scalable and embodiment-rich data source. A multi-stage curation pipeline ensures signal quality across all heterogeneous sources.
Beyond manipulation data, a curated vision-language mixture including novel embodied chain-of-thought and egocentric video understanding data preserves the VLM backbone's perceptual and reasoning capabilities during VLA training.

\textbf{A new standard for evaluating VLA generalization.} We introduce OOD evaluation settings including LIBERO-Plus, RoboTwin-Clean2Rand, RoboCasa365, and EBench, alongside RoboTwin-IF, a benchmark that diagnoses genuine language conditioning, and RoboTwin-XE, a benchmark for zero-shot cross-embodiment transfer. We argue that in-domain metrics are insufficient proxies for foundation model capability and that OOD transfer is the correct measure. \ours substantially outperforms existing VLA models across all OOD settings while achieving state-of-the-art on standard benchmarks.

\textbf{Real-robot validation across diverse deployment scenarios.} 
We validate \ours on four physical platforms (AgileX ALOHA, Franka, UR, and ARX) across in-domain, out-of-domain, few-shot adaptation, and zero-shot cross-embodiment transfer settings.
On the RoboChallenge Table30 v1 generalist track, \ours ranks 1st.
Results confirm that the generalization capabilities of \ours hold under real-world deployment conditions.
% We validate on physical platforms spanning fine-grained dexterous tasks such as cap screwing and cross-embodiment transfer scenarios, confirming that the generalization gains observed in simulation benchmarks carry over to real-world deployment.

%%%%%%%%%%%%%%%%%%%%%%%%%%%%%%%%%%%%%%%%%%%%%%%%%

% Achieving this standard requires confronting two interconnected challenges. First, the heterogeneity of state, action, and end-effector spaces across robot platforms makes it difficult to train a single model on multi-source data while preserving cross-embodiment structure, demanding careful alignment at the representation, motion, and behavioral levels. Second, existing robot datasets are dominated by costly teleoperated demonstrations concentrated in narrow manipulation setups, leaving the vast reservoir of naturally available egocentric human manipulation video largely untapped and limiting the diversity and scale of training data.

% These two challenges point naturally to our approach. To address representational heterogeneity, \ours introduces a unified alignment framework operating at three levels. At the representation level, an 80-dimensional canonical vector accommodates diverse robot morphologies within a single template, with a per-dimension binary mask ensuring gradients flow only through semantically populated dimensions. At the motion level, a camera-frame delta pose parameterization grounds end-effector actions directly in the visual domain, ensuring that visually similar motions are numerically proximate across embodiments and coordinate frames. At the behavioral level, an in-context policy adaptation mechanism reads intra-episode execution history as an implicit embodiment identifier, enabling the model to adjust its behavior based on the kinematic and interaction patterns of the current episode without any parameter update. To address data scarcity and diversity, \ours combines a rigorous multi-stage data curation and pre-processing pipeline with a human-to-robot synthesis pipeline that converts egocentric hand demonstrations into robot trajectories across 15 platforms, bringing the total training corpus to $\sim$30,000 hours. Throughout training, a dual-stream co-training strategy jointly optimizes over manipulation data and a curated vision-language stream, preventing the VLM backbone's pretrained perceptual and reasoning capabilities from eroding under action prediction optimization.

% In this report, we introduce \textbf{\ours}, a genuinely generalizable Vision-Language-Action foundation model for robotic manipulation built upon Qwen-VL. Our contributions are:

% \textbf{Scalable data infrastructure.} We consolidate $\sim$30,000 hours of manipulation data from open-source robot datasets and $\sim$1.8M egocentric human demonstrations. A human-to-robot synthesis pipeline renders any of 15 robot embodiments from a single egocentric demonstration, providing a scalable and inherently embodiment-agnostic data source. A multi-stage curation and pre-processing pipeline ensures signal quality across all heterogeneous sources.

% \textbf{Unified multi-level alignment.} We address representational heterogeneity through three complementary mechanisms: a canonical state-action vector with per-dimension binary masking, a camera-frame delta pose parameterization that grounds end-effector actions in the visual domain, and an in-context policy adaptation mechanism that reads intra-episode execution history as an implicit embodiment identifier. To prevent action prediction optimization from eroding the pretrained perceptual and reasoning capabilities that underpin generalization, training is conducted under a dual-stream co-training strategy that jointly optimizes over manipulation data and a curated vision-language stream.

% \textbf{Strong generalization and a new instruction-following benchmark.} \ours achieves state-of-the-art performance on LIBERO, RoboCasa-GR1, and RoboTwin, and demonstrates substantially stronger OOD generalization than existing VLA models across LIBERO-Plus, RoboTwin-Easy2Hard, RoboCasa365, and EBench. We further introduce a new instruction-following benchmark that tests genuine language conditioning, and validate on real-robot platforms spanning dexterous tasks and cross-embodiment transfer.

% The remainder of this report is organized as follows. Section~\ref{sec:data} describes our data sources and curation pipeline. Section~\ref{sec:model} presents the model architecture. Section~\ref{sec:training} details the training and post-training strategy. Section~\ref{sec:system} describes the agentic and inference systems. Section~\ref{sec:experiments} presents experimental results and ablations. We conclude in Section~\ref{sec:conclusion}.

%%%%%%%%%%%%%%%%%%%%%%%%%%%%%%%%%%%%%%%%%%%%%%%%%

% Despite rapid progress, however, the generalization of existing VLA models is more limited than it appears. Models such as OpenVLA~\citep{}, GR00T~\citep{bjorck2025grootn1}, and RoboVLMs~\citep{} do report out-of-distribution evaluations, yet these typically involve minor visual perturbations, \textit{e.g.} swapping object instances or backgrounds, while keeping the same robot embodiment, task structure, and workspace layout as training. True OOD generalization across embodiments, task categories, or substantially different environments remains largely out of reach. In practice, deploying such models on a new platform or workspace still demands extensive re-collection of demonstrations and retraining from scratch. This is not merely an inconvenience. It is a diagnostic signal. On several standard manipulation benchmarks, fine-tuning a pretrained VLA model on the provided training split yields performance comparable to training from scratch on the same data~\citep{}, suggesting that the pretrained prior contributes little transferable value. A foundation model that offers no advantage over random initialization at transfer time has not learned a genuinely generalizable representation of manipulation.

% The robotics community has implicitly recognized this gap. $\pi_{0.5}$~\cite{} is perhaps the clearest illustration: its widespread adoption stems not from topping standard benchmarks, but from a simple practical virtue. Users can fine-tune it on a modest amount of data collected on their own robot and obtain strong, reliable performance on their target task. This ease of adaptation is what makes it genuinely useful in the field. It suggests that the correct criterion for a VLA foundation model is not in-distribution task success, but downstream adaptability, namely how efficiently a user can transfer the model to a new embodiment and task with minimal domain-specific data.

% Achieving this standard requires confronting three interconnected limitations of prior work. First, existing robot datasets are dominated by costly teleoperated demonstrations concentrated in narrow laboratory setups, leaving the vast reservoir of naturally available egocentric human manipulation video largely untapped, a missed opportunity for scaling both data diversity and embodiment coverage. Second, the heterogeneity of state, action, and end-effector spaces across robot platforms has prevented prior approaches from learning a unified physical representation. Models either train separate heads per embodiment, forfeiting cross-embodiment synergy, or adopt overly compressed representations that discard the kinematic structure needed for precise control. Third, existing VLA models are stateless within each episode, making no use of the robot's own historical state--action trajectory, a structured behavioral context that is freely available at deployment and directly informative about how the robot has been interacting with its environment.

% These three gaps point naturally to our approach. We address data scarcity through a scalable human-to-robot synthesis pipeline that converts egocentric hand demonstrations into multi-embodiment robot training data, a source that is naturally cross-embodiment by construction. We address representational fragmentation through a unified canonical state-action space paired with a camera-frame end-effector representation, ensuring that actions which appear visually similar are numerically proximate across embodiments and enabling genuine cross-embodiment knowledge transfer. We address episodic statelessness through an in-context policy adaptation mechanism that conditions action prediction on a rolling window of intra-episode historical chunks, serving as an implicit embodiment identifier that encodes how this particular robot has been executing the task. And we show that strong OOD generalization, once achieved, unlocks a qualitatively new capability. Through hierarchical agentic composition, the model can generalize to novel objects, locations, and task combinations that were never seen during training.

% \textbf{Unified cross-embodiment representation.} We design a canonical 80-dimensional state-action vector structured around semantic groups including joint angles, end-effector pose, gripper state, and dexterous hand joints, accommodating single-arm grippers, dual-arm systems, humanoids, and mobile manipulators within a single template. A per-dimension binary mask restricts training loss to semantically populated dimensions. Complementing this, we adopt a camera-frame delta pose representation for end-effector motion, ensuring that visually similar actions are numerically proximate across embodiments and grounding the action space directly in the visual domain. Together, these design choices maximize cross-embodiment synergy without sacrificing embodiment-specific structure.

% \textbf{Scalable data infrastructure with human-to-robot synthesis.} We consolidate a large-scale corpus spanning open-source robot datasets and $\sim$1.8M egocentric human manipulation trajectories. An end-to-end synthesis pipeline converts MANO hand keypoints into robot end-effector trajectories, inpaints human arms using SAM3 and ProPainter, and renders any of 12 robot embodiments via a MuJoCo digital twin, allowing a single human demonstration to yield training data across multiple robots. Because egocentric human data is inherently embodiment-agnostic, this pipeline is a natural source for cross-embodiment generalization at scale. A five-stage data curation pipeline ensures signal quality across all heterogeneous sources.

% \textbf{In-context policy adaptation.} \ours conditions current action prediction on a rolling window of historical observation--state--action chunks from the same episode, serving as an implicit embodiment identifier that encodes the kinematic signature and behavioral style of the current deployment. Each context chunk mirrors exactly the format of the model's own inference procedure, requiring no architectural modification. Stochastic context sampling during training, which randomizes the number and temporal position of history chunks, prevents shortcut learning and forces the model to extract consistent behavioral style from any intra-episode subset.

% \textbf{Agentic VLA system.} The strong OOD generalization of \ours is not merely a benchmark metric. It enables a qualitatively new mode of deployment. We build a hierarchical agentic system in which a high-level planner decomposes long-horizon tasks into sequences of subtask instructions, each grounded and executed by \ours as a low-level policy. Because \ours can reliably act on novel objects and unseen spatial configurations, the system achieves compositional generalization: tasks and objects unseen during training can be handled through instruction composition, without any policy retraining.

% \textbf{State-of-the-art performance and genuine OOD generalization.} \ours achieves state-of-the-art results on standard fixed-table manipulation benchmarks, including fine-grained dexterous tasks such as cap screwing that demand sub-centimeter precision. More importantly, it demonstrates substantially stronger OOD generalization than existing VLA models including $\pi_{0.5}$, across novel objects, unseen scenes, and cross-embodiment transfer. When fine-tuned on a small amount of domain-specific data, the strong prior provided by \ours translates directly into faster convergence and higher final performance, which is the standard by which a foundation model should ultimately be judged.
